\documentclass{ximera}
\title{Elastic Bands, Jar Lids and Measuring the Earth}

\newcommand{\pskip}{\vskip 0.1 in}

\begin{document}
\begin{abstract}
A few related problems.
\end{abstract}
\maketitle

\section{The Jar Lid and an Elastic Band}
\begin{question} \label{4Ef4tDE}

You wrap an elastic band around a jar lid of radius $a$ cm and stretch the band as shown below.

\begin{onlineOnly}
    \begin{center}
\desmos{ymesxpazea}{900}{600}
\end{center}
\end{onlineOnly}

\href{https://www.desmos.com/calculator/ymesxpazea}{Jar Lid 44}

\begin{enumerate}

\item Drag the point $P$ in the worksheet above to see how the elastic band changes. Pay particular attention to the points of tangency $Q$ and $R$ and describe what you see.

\begin{freeResponse}
\end{freeResponse}

\item Sketch by hand a graph of the function $L=f(h)$ that express the length of the band (measured in cm) in terms of $a$ and the distance $h$ between the point $P$ and the center of the lid, also measured in cm. Pay particular attention to oblique asymptotes.

\item Find an expression for the function $L=f(h)$. include the domain. Do this as follows.

\item Check that your expression has the correct units.
\begin{freeResponse}
\end{freeResponse}

\begin{enumerate}
\item First express the length of one straight part $\overline{PQ}$ of the elastic band in terms of $h$ and $a$. 

The length is 
\[
  PQ = \answer{\sqrt{h^2-a^2}}.
\]

\item Next find an expression for the angle $\theta = \angle QOP$.

The angle $\theta$ is
\[
  \theta = \answer{\arccos(a/h)}.
\]

\item Then find the length of portion of the band wrapped around the lid.

This length is 
\begin{align*}
\text{long arc} QR &= 2\pi a - \answer{2a\arccos(a/h)} \\
                           &= 2a \arccos (\answer{-a/h})
\end{align*} 

\item So the length of the band in cm is
\[
 f(h) = \answer{2\sqrt{h^2 - a^2}+2a\arccos(-a/h)} \, , \, h\geq \answer{a} .
\]
\end{enumerate}
\end{enumerate}
\end{question}

\section{Measuring the Radius of the Earth}

\begin{question} \label{Q6:RightTriangle6}
On the vernal equinox at a location on the equator, the sun rises due east at 6am, passes directly overhead at noon, and sets due west at 6pm. 

Imagine you and a friend are at the beach watching the sunset from a point on the equator on the vernal equinox. You lie on the beach and your friend stands on the balcony of your hotel, $h$ feet above the ground. As illustrated below, your friend will see the sunset later than you. 


\begin{onlineOnly}
    \begin{center}
\desmos{y8ijqldt5y}{900}{600}
\end{center}
\end{onlineOnly}

\href{https://www.desmos.com/calculator/y8ijqldt5y}{142: Skyscraper Sunset 4}

\begin{enumerate}
\item Find a function
\[
   t = f(h) \, , h\geq 0 ,
\]
that takes as an input the height $h$ of your friend above the ground (measured in meters) and returns as an output the difference in time (measured in seconds) between the moments the two of you see the sun dip below the horizon. Take the radius of the earth to be $a$ kilometers.

\item How much later does your friend see the sun set if she is $100$ meters above the ground?

\item Find a function 
\[
   h = f^{-1}(t)
\]
that express the height of your friend above the ground in terms of the time difference between the observed sunsets. %This suggests a way to approximate the radius of the earth. How?

\item How high above the ground is your friend if she sees the sun set two minutes later than you do?

\item Find a function
\[
     a = g(h,t)
\]
that expresses the radius of the earth (in km) in terms of your friend's height above the ground (in meters) and the difference $t$ in the times of the observed sunsets (measured in seconds).
\end{enumerate}



\end{question}

\end{document}